\section{Experiments}
\label{sec:exp}

\subsection{Setup}
We use the public 118k-polymer benchmark with 42\,M nodes. Evaluation follows
the scaffold split of \citet{wang2020scaffold}: five folds, monomers grouped by
scaffold so that no scaffold appears in both training and test. All models
share the same budget: 3.7\,M parameters, 200 epochs, batch size 64, AdamW with
learning rate $1\times10^{-3}$ and weight decay $1\times10^{-2}$. Reported
numbers are means over five folds; variance is in the released logs.

\subsection{Main results}
Table~\ref{tab:metrics} compares SAGN with two baselines at matched parameter
count. SAGN attains the lowest error on both metrics; the gap over GIN is
largest on the dielectric subset, where long-range terms dominate.

\begin{table}[htbp]
\centering
\caption{Property prediction error on the 118k-polymer scaffold split
(mean over five folds). Lower is better for both metrics.}
\label{tab:main-res}
\begin{tabular}{lllllll}
\toprule
Model & Backbone & Pretraining corpus & Split protocol & MAE (K) & RMSE (K) & Params (M) \\
\midrule
SAGN (ours) & GINE + sparse attention, 6 layers & 118k polymers, 42M nodes, 4 bond types & scaffold split, 5 folds, seed 42 & 6.82 & 11.40 & 3.7 \\
GIN baseline & GIN, 5 layers, sum readout, 300 hidden & 118k polymers, 42M nodes, 4 bond types & scaffold split, 5 folds, seed 42 & 8.15 & 13.66 & 3.1 \\
MPNN baseline & edge-conditioned MPNN, 4 layers & 118k polymers, 42M nodes, 4 bond types & scaffold split, 5 folds, seed 42 & 8.44 & 14.02 & 3.4 \\
\bottomrule
\end{tabular}
\end{table}

\subsection{Ablation}
Removing the low-rank branch costs $0.91$\,K MAE; removing sparsification
($k=|V|$) costs $0.34$\,K but triples wall-clock per epoch. Both variants stay
inside the same parameter budget, so the comparison is budget-matched.
